\section{Propriétés des opérateurs discrets} \label{sec:properties_of_the_operators}

\subsection{Formules de Green discrètes}

Les opérateurs discrets introduits sont en dualité. En effet, grâce aux produits scalaires dans la \Cref{sec:produits_scalaires}, on peut écrire des formules de Green discrètes. C'est cette propriété de dualité qui a donné son nom à la méthode \ddfv.

\begin{theo}[Formules de Green discrètes] \label{theo:discrete_green_formulae}
Pour tout $\xi^\ensdiamantplat, u^\ensddfvplat \in (\R^2)^\ensdiamantplat \times \R^\ensddfvplat$, on a
\[
\prodscalddfv{\div^\ensddfvplat \xi^\ensdiamantplat}{u^\ensddfvplat} = - \prodscalD{\xi^\ensdiamantplat}{\grad^\ensdiamantplat u^\ensddfvplat}
\quad
\text{et}
\quad
\prodscalddfv{\scarot^\ensddfvplat \xi^\ensdiamantplat}{u^\ensddfvplat} = \prodscalD{\xi^\ensdiamantplat}{\rot^\ensdiamantplat u^\ensddfvplat}.
\]
\end{theo}
\begin{demo}[{\cite[Prop. 2.3]{domelevo_finite_2005}}]
\areprendre{
Let $\xi^\ensdiamant$ be a discrete vector field and $\grad^\ensdiamant u^\ensddfv$ the discrete gradient of a function $u^\ensddfv$, respectively defined by the values $\xi^{\diamant{D}}$ and formula \ref{eq:gradient_discret} on each of diamond-cells $\diamant{D} \in \ensdiamant$. There holds:
\begin{equation}
    \prodscalD{\xi^\ensdiamant}{\grad^\ensdiamant u^\ensddfv} = \sum_{\diamant{D} \in \ensdiamant} \mesure(\diamant{D}) \, \xi^{\diamant{D}} \cdot \grad^{\diamant{D}} u^\ensddfv = 
    \sum_{\diamant{D} \in \ensdiamant} \frac{1}{2} \, \xi^{\diamant{D}} \cdot \mathopen{}\Big[ (u_{\primal{L}} - u_{\primal{K}}) \mesure(\sigma) \, \vect{\nu} + (u_{\mailleduale{L}} - u_{\mailleduale{K}}) \mesure(\dual{\sigma}) \, \dual{\vect{\nu}} \Big].
\end{equation}
\begin{equation}
    \prodscalD{\xi^\ensdiamant}{\grad^\ensdiamant u^\ensddfv} = - \frac{1}{2} \sum_{\primal{K} \in \ensprimal} u_{\primal{K}} \mathopen{}\left( \sum_{\diamant{D}_{\sigma,\dual{\sigma}} \in \ensdiamant_{\primal{K}}} \mesure(\sigma) \, \xi^{\diamant{D}} \cdot \vect{\nu} \right)
    - \frac{1}{2} \sum_{\mailleduale{K} \in \ensdual} u_{\mailleduale{K}} \mathopen{}\left( \sum_{\diamant{D}_{\sigma,\dual{\sigma}} \in \ensdiamant_{\mailleduale{K}}} \mesure(\dual{\sigma}) \, \xi^{\diamant{D}} \cdot \dual{\vect{\nu}} \right).
\end{equation}
}
\end{demo}

\subsection{Compositions des opérateurs discrets} \label{sec:compositions_discrete_operators}

\cite[Sec. 4.2]{delcourte_discrete_2007} The aim of this section is to {\color{red}verify} a discrete analogue of the following continuous identities: $\div \rot = 0$, $\scarot \grad = 0$ and $\scarot \rot = - \div \grad$\ .

\begin{prop}
Pour tout $u^\ensddfvplat \in \R^\ensddfvplat$, on a
\begin{align}
    \div^\ensddfvplat (\rot^\ensdiamantplat u^\ensddfvplat) &= 0 \\
    \scarot^\ensddfvplat (\grad^\ensdiamantplat u^\ensddfvplat) &= 0. 
\end{align}
\end{prop}

\begin{prop}
Pour tout $u^\ensddfvplat \in \R^\ensddfvplat$, l'égalité suivante est vérifiée
\[
\scarot^\ensddfvplat \rot^\ensdiamantplat u^\ensddfvplat = - \div^\ensddfvplat \grad^\ensdiamantplat u^\ensddfvplat.
\]
\end{prop}

\subsection{Discrétisation de la décomposition de Helmholtz--Hodge}

\cite[Sec. 4.3]{delcourte_discrete_2007} In the continuous case, the Hodge decomposition for non simply connected domains reads:
\begin{equation} \label{eq:413}
(\L^2)^2 = \grad V \oplusperp \rot W,
\end{equation}
with $V = \ens[\big]{\phi \in \H^1 \tq \int_\Omega \phi = 0}$ and $W = \ens[\big]{\psi \in \H^1 \tq \cdot}$. To prove an analogous property in the discrete case, we rely on the following result:

\begin{lemme}[Euler's formula]
For a non simply connected \dots
\[
\chi = \dots
\]
\end{lemme}

\begin{lemme}[Identité d'Euler{\cite{gregoire_decomposition_2025}}]
Soient $\chi$ la caractéristique d'Euler d'une surface donnée, un graphe sur cette surface avec $A = \#{\text{arêtes}}$, $S = \#{\text{sommets}}$ et $F = \#{\text{faces}}$. Alors
\[
S - A + F = \chi.
\]
\end{lemme}

\begin{corol}[{\cite{gregoire_decomposition_2025}}]
\begin{itemize}
    \item La dimension de l'espace vectoriel des fonctions scalaires sera : $S + F$.
    \item La dimension de l'espace vectoriel des fonction vectorielles sera : $A$.
    \item Si $\chi$ est nul, ce qui est le cas pour le tore, nous serons assurés d'avoir des opérateurs tel que le gradient ou la divergence par exemple carrés. Dans le cas contraire, la cohomologie devra intervenir dans la décomposition.
\end{itemize}
\end{corol}

On peut maintenant énoncer la décomposition de Helmholtz--Hodge discrète.
\begin{theo}[Décomposition de Helmholtz--Hodge discrète]
Soit $\xi^\ensdiamantplat \in (\R^2)^\ensdiamantplat$ un champ de vecteurs discret défini par ses valeurs sur les diamants plats. Il existe un unique couple $(\phi^\ensddfvplat, \psi^\ensddfvplat) \in \Rddfvmn \times \R^\ensddfvplat$ tel que
\begin{equation} \label{eq:415}
    \xi^\ensdiamantplat = \grad^\ensdiamantplat \phi^\ensddfvplat + \rot^\ensdiamantplat \psi^\ensddfvplat.
\end{equation}
De plus, la décomposition \ref{eq:415} est orthogonale. 
\end{theo}

\begin{demo}
    
\end{demo}

\begin{remarque}
La condition $\phi^\ensddfvplat \in \Rddfvmn$ est un analogue discret de la condition $\int_\Omega \phi = 0$ qui apparaît dans la définition de l'espace $V$ dans \ref{eq:413}. 
\end{remarque}

\cite[Sec. 3.2.5]{gregoire_decomposition_2025} Nous allons aussi définir la divergence et la divergence perp. Pour cela, nous allons utiliser la dualité grâce à la formule d'intégration par parties
\[
\prodscalddfv{\div^\ensddfv \xi^\ensdiamant}{u^\ensddfv} = - \prodscal{\xi^\ensdiamant}{\grad^\ensdiamant u^\ensddfv}_\ensdiamant.
\]
L'objectif est de résoudre
\begin{equation} \label{eq:DHH}
\xi^\ensdiamant = \grad^\ensdiamant \varphi^\ensddfv + \rot^\ensdiamant \psi^\ensddfv + r^\ensdiamant.
\end{equation}
En prenant la divergence, on obtient
\[
\div^\ensddfv \xi^\ensdiamant = \laplacien \varphi^\ensddfv.
\]
On ne va jamais calculer explicitement la divergence et la divergence perp, mais on sera amener à résoudre l'équation
\[
\forall v^\ensddfv \in \R^\ensddfv, \quad
\prodscalD{\xi^\ensdiamant}{\grad^\ensdiamant v^\ensddfv} = \prodscalD{\grad^\ensdiamant \varphi^\ensddfv}{\grad^\ensdiamant v^\ensddfv}.
\]
Pour obtenir l'équation en $\psi^\ensddfv$, on applique $\scarot^\ensddfv$ à l'équation \ref{eq:DHH}:
\[
\cdots
\]
\begin{remarque}
La matrice $\rotmat^\top_\ensdiamant \masse_\ensdiamant$ est symétrique positive mais pas définie. Néanmoins, cela n'est pas gênant car la méthode du gradient conjugué converge à un vecteur du noyau de $\gradmat_\ensdiamant$ près. Or le noyau de $\gradmat_\ensdiamant$ contient l'homologie ainsi que des modes hyper-oscillants. De plus, la composante dans le noyau de $\gradmat$ est préservée dans toutes les étapes du gradient conjugué et en prenant comme solution initiale un vecteur nul, cela nous assure de n'obtenir aucun de ces modes tout en convergeant. 
\end{remarque}